% =========================================================================
% conclusion.tex
% =========================================================================

\section{Conclusion and Future Work}
\label{sec:conclusion}

I have presented a systematic empirical study of KV cache vector
quantization at attention head dimension $\dhead = 64$, a setting not
evaluated in the four source papers built upon.
The experiments on GPT-2 Medium identify two failure modes at this scale: high QJL residual variance and the collapse of FibQuant angular codebooks. I address these through a simple fallback heuristic that disables QJL residual estimation below $\dhead = 128$, as well as through MC-TurboQuant, which combines WHT vector quantization with Gated Residual Memory.

Disabling the QJL residual calculation raises nearest-neighbor Recall@1 from 38.8\% to 67.4\% at $3.2{\times}$ compression, without increasing computational cost.
MC-TurboQuant maintains 97.5\% average NIAH recall under quantization,
demonstrating that recurrent memory caching and vector quantization
compose without quality degradation.
Empirical distortion remains within $1.83$--$2.46{\times}$ of the
Shannon lower bound across two to five bits, confirming that the
theoretical optimality of rotation-based scalar quantizers extends
to $\dhead = 64$.

The immediate next step is cross-dimension generalization: testing
Adaptive QJL at $\dhead = 128$ on a larger model to determine whether
the QJL residual correction is beneficial at that scale (as the
source papers suggest) or whether the Adaptive approach provides
improvements there as well.
This experiment was not completed in the present study.
Additional directions include implementing fused CUDA kernels to
realize the theoretical latency advantages of the WHT, scaling the
MC-RNN training experiment to match the parameter and data budgets
of the original Memory Caching paper, and evaluating on models with
intermediate head dimensions (\eg, $\dhead = 96$) to locate the
precise threshold at which QJL residual correction transitions from
harmful to beneficial.
